\chapter{电缆机器人的路径规划算法}\label{chap:path_planning}

\section{路径规划目标}
电缆机器人主要工作于高空输电线路等复杂环境中，其路径规划的核心目标是在确保机器人自身运动稳定与安全的前提下，寻找一条从当前起点到终点的无碰撞、且在时间和能量上达到最优的可行路径。具体而言，针对本文的五连杆轮腿式电缆机器人，路径规划目标可归纳为以下几个方面：

1. \textbf{环境适应与实时避障}：高压输电线上通常存在防震锤、绝缘子串、悬垂线夹等多种类型的金具障碍物。规划算法需根据传感器信息识别障碍位置，并自主决定跨越或是匀速绕行；
2. \textbf{姿态与动态稳定性保证}：受到高空风偏和线缆弧垂扰动的影响，结合前文建立的动力学模型，机器人在规划路径及越障过程中的质心波动、机体俯仰角必须被限制在安全范围内，防止发生脱线或翻转侧翻；
3. \textbf{多目标效能优化}：为了尽最大限度提高高空作业的续航能力，规划出的路径速度与虚拟腿长曲线应尽量平缓，减少电机的剧烈加减速与冗余输出做功，实现能量的高效利用。

\section{传统路径规划算法}
在移动机器人领域，经典的传统路径规划算法主要包括基于图搜索的算法（如 A*、Dijkstra）、基于采样的算法（如 RRT 快速随机搜索树）以及人工势场法（APF）、局部避障算法 DWA 等。这里主要介绍两种典型的方法及其在电缆机器人上的适用性考量。

\subsection{A* 算法}
A* 算法是一种结合了全局启发式信息的图搜索策略。它通过维护一张搜索代价值和预估距离目标点代价的启发函数列表，在已知的离散化栅格地图中迅速收敛寻找最短路径。
对于电缆机器人而言，在全局范围内提前获知电缆线上的故障点及防震锤位置时，使用 A* 算法可以在宏观层面快速且无偏地生成避障及检修跨越的“站位点”序列。然而，当电缆受到风力发生动态形变时，A* 算法生成的静态理想路线难以直接适配由于悬垂带来的高度与重力势能变化，必须依赖频繁重规划。

\subsection{DWA 动态窗口法}
DWA（Dynamic Window Approach）是一种常用且高效的局部路径规划和避障算法。它主要考虑了机器人的运动学约束（最大速度、极大加速度约束），在机器人当前的速度空间（$v, \omega$）内采样一个满足刹车距离安全的“速度窗口”，并根据目标导向、避障间距和速度期望值计算最优的一组控制速度以驱动机器人运动。
将 DWA 用于电缆机器人时，可以将电缆一维行走与五杆虚拟腿长的联合控制构成多维状态空间的避障。然而，该算法虽然计算高效，但在避开密集金具序列且动态非线性强耦合的控制情境下（如机器人质心高度的剧烈被动改变）容易陷入局部极小值停滞不前，难以满足高度复杂的实时跨越要求。

\section{强化学习路径规划算法}
针对传统基于模型与搜索算法在高维、非线性连续状态空间控制及未知动态环境中的不足，基于深度强化学习（Deep Reinforcement Learning, DRL）的路径规划为电缆机器人提供了更为先进的智能化解决思路。强化学习算法通过智能体（Agent）与环境的不断试错交互来学习最佳策略，摆脱了对精确全局地图和复杂运动学求逆的高度依赖，展现出了绝佳的泛化能力。

在具体的电缆机器人路径规划任务中，常用的连续动作空间算法包括 DDPG（深度确定性策略梯度算法）、PPO（近端策略优化算法）以及 SAC（柔性动作评价算法）等。以该五杆轮腿机器人的电缆越障及行进规划为例，其强化学习模型可进行如下抽象定义：
\begin{itemize}
    \item \textbf{状态空间 (State)}：由本体感知与环境感知融合构成。本体感知涵盖机器人当前的俯仰姿态（IMU 数据）、各关节角（$\phi_1, \phi_4$）及角速度等；环境感知则利用视觉或雷达传感获取正前方最近金具或障碍物的距离特征。
    \item \textbf{动作空间 (Action)}：由强化学习策略网络（Actor）输出的连续控制动作，可直接解耦并映射为轮毂电机的纵向期望速度增量，以及结合 VMC 模型所需的虚拟腿长 $L_0$ 和虚拟摆角 $\varPhi_0$ 操作值，从而调节机器人整体位姿。
    \item \textbf{奖励函数 (Reward)}：引导智能体进化的主要依据。其应包含对向目标点靠近并抵达的终极奖励、对产生摩擦碰撞及姿态失稳（例如翻转、力矩超限）的严厉惩罚。同时可加入控制量平滑变化的惩罚项，促使智能体策略更加节能、稳健。
\end{itemize}

经过在仿真的高空电缆环境中的充分训练验证，电缆机器人便可自主获得集环境动态避障、车体位姿平衡与底层 VMC 控制命令输出为一体的“端到端”最优路径规划模型。


